\chapter{Introduction}

Compared to written language, speech is a complex, multi-faceted dimensional signal. Linguistic analysis distinguishes between two primary components: the phonetic segments or phonemes on the one hand, and prosodic features like pitch movement, amplitude, and duration on the other. Standard orthographies of written language capture mostly only the phonemic part of speech, which can be reduced to discrete elements that are represented by letters or graphemes. To be sure, written language has developed its special devices, such as headings, paragraphs, or bullet lists. But compared to the human voice, orthographic representation of text is "flat", hardly able to capture the expressive dimension of language \cite{pataca2021speech}. This straight-jacket is felt in current digital communication, where users frequently attempt to get rid of these limitations by additional symbols, like emojis \cite{zappavigna2024}, but also by non-standard orthography, using "visual prosody" like all-caps to signal prosodic highlighting, emotive arousal, or additional iconic meanings, such as by vowel elongation (e.g., \textit{booooring}) to mimic intonational lengthening \cite{fuchs2019}, \cite{heath2020no}.

With the advent of text-to-speech and speech-to-text systems, the gap between speech and writing has narrowed dramatically. However, automated speech recognition (ASR) prioritizes orthographic transcription over fine-grained acoustic detail, as represented, e.g., by the International Phonetic Alphabet (IPA) \cite{internationalphoneticassociation1999ipa}. This creates a gap between the richly expressive nature of the human voice and the simplified representation provided by standard orthography. It may be the optimal way to go for applications like dictation devices for business, but it may be suboptimal, for example, for subtitle generation for movies.

In this thesis, we will conceive, and partially develop, a tool, called SpeechPrint, which enables the integration of prosodic features in the written representation of spoken language. The range of applications of such a tool is considerable: from lively subtitle representations over voice input to create more expressive chats to the artistic visualization of poem recitals. We will discuss such possible uses, but for the purpose of this thesis, we will focus on the specific task to create a symbolic representation of prosodic features that is able to support such uses. To be specific, we will develop a tool to assist with the transcription of speech for linguistic research. This rather specific use scenario helps us focus on extracting of relevant prosodic features, in particular lengthening, pitch movement, and amplitude. We think that it should also be a useful device in fields like conversation analysis and field work on under-resourced languages.

There are excellent tools available to visualize and study speech. One, Praat, was initiated more than 20 years ago, has been continuously developed, and is now used all over the world in linguistic and phonetic labs, as well as for the study of animal communication. But with the advent of ASR, there are increasingly capable tools that assist with the arduous tasks of the transcription and annotation of speech. They have to be trained for the languages they are used for, which means they are better at large, widely used languages of economic importance than at smaller languages with fewer resources. However, they typically aim for standard orthographic representations that are time-aligned at the level of utterances. Programs that time-align speech segments, such as phonemes, are rarer, and those that also attempt to render prosodic features are highly specialized.

Our goal is to investigate and test available modules with this more specialized orientation and bundle them into a single device that provides time-aligned phonemes and a prosodic representation. The immediate goal is to create a tool that should be of help to linguists who are not specialized in phonetic research, such as conversation analysis \cite{hutchby2008} and linguistic field work \cite{thieberger2012}, \cite{rozhanskiy2021}. SpeechPrint, at the planned stage, should provide a pre-transcription analysis with phonemes and time-aligned syllables, and a symbolic representation of pitch and amplitude matched to the syllables. This pre-analysis can be used to help with a more precise, human-guided transcription effort. The result can be only as good as the current tools that we have found in the current stage of development. It is not surprising that they are better at large languages like English and Spanish than at smaller ones.

The thesis proceeds as follows. In Chapter \ref{StateOfArt}, we will give a very brief sketch of the current state of spoken language research, as far as it is relevant for this thesis. It focuses on phonemes and prosody, and introduces two main tools, Praat and ELAN. In Chapter \ref{ASR}, we will discuss Automated Speech Recognition and important recent tools that have been developed from this research, in particular phoneme aligners and prosody analyzers. Chapter \ref{speechprint} is devoted to the description of the SpeechPrint tool, and in Chapter \ref{testing} we will put it to various tests. The final chapter, \ref{scenarios}, consists of a discussion of what is achieved, of what remains to be achieved, and, in particular, of some usage scenarios of representing prosodic features typographically, or more generally, graphically, to bring alive the vividness and expressiveness of speech in the visual domain.

##  CHECK THIS FOR CONSISTENCY

